\documentclass[11pt]{article}
\usepackage[margin=1in]{geometry}
\usepackage[T1]{fontenc}
\usepackage[utf8]{inputenc}
\usepackage{lmodern,microtype,amsmath,amssymb,amsthm,mathtools}
\usepackage[colorlinks=true,linkcolor=blue,citecolor=blue,urlcolor=blue]{hyperref}
\hypersetup{pdftitle={Exponential online mistakes at every width after arbitrarily small Gaussian anisotropy},pdfauthor={ChatGPT-6 Astra; Kaarel Hanni; Dmitry Vaintrob}}
\newtheorem{theorem}{Theorem}
\newtheorem{lemma}[theorem]{Lemma}
\newtheorem{proposition}[theorem]{Proposition}
\newtheorem{corollary}[theorem]{Corollary}
\newtheorem{remark}[theorem]{Remark}
\newcommand{\R}{\mathbb R}
\newcommand{\E}{\mathbb E}
\newcommand{\Pp}{\mathbb P}
\newcommand{\sgn}{\operatorname{sign}}
\newcommand{\ess}{\operatorname{ess}}
\newcommand{\wt}{\operatorname{wt}}
\title{Exponential online mistakes at every width\protect\\after arbitrarily small Gaussian anisotropy}
\author{ChatGPT-6 Astra \and Kaarel H\"anni \and Dmitry Vaintrob\thanks{Authors are listed in alphabetical order.}}
\date{6 November 2026}
\begin{document}
\maketitle

\begin{abstract}
We prove an exponential online-mistake lower bound for a Bayesian tanh
network at every hidden width. The target is ordinary majority on $d$ odd
input bits: it has a linear-size Boolean circuit and a width-one teacher
with every weight at most one. The learner uses one hidden layer, a
constant input coordinate, a linear readout, global absolute output
margin greater than $1/2$, and exact training labels.
The first-layer Gaussian variances have an arbitrarily small, explicitly
specified coordinate asymmetry. One fixed input permutation forces at least
$2^{\lfloor d/8\rfloor}$ mistakes in the strong-prior limit, and forces all
of them with arbitrarily high probability at sufficiently small positive
variance. The covariance perturbation and input order are independent of
the learner width. The proof bounds the energy of every nonlinear network
using its essential input coordinates; it does not restrict the posterior
to scalar subnetworks or choose an unfavorable classifier from a tie.
For $d=9q$ with odd $q$, a second construction uses only $2q$ calibration
examples and forces $9^{(q+1)/2}$ mistakes, with an explicit
inverse-polynomial upper bound on the anisotropy amplitude.
The same priors also have width-independent polynomial IID
sample complexity for this target. Neither construction establishes
the same lower bound for the exactly isotropic Gaussian prior: at exact
isotropy, the relevant candidate tier has equal masses and only a
logarithmic cumulative sampling-error bound before it is exhausted.
\end{abstract}

\section{Model and result}

An earlier width-one anisotropic construction appears
in~\cite{and-gaussian}. The essential-coordinate bound below permits a
comparison with every wide nonlinear competitor. It also yields a
linear-calibration construction and a polynomial IID guarantee for
the same prior and target.

Fix an odd input dimension $d\ge17$ and a hidden width $N\ge1$.
Inputs lie in $\{-1,1\}^d$. Write $\bar x=(1,x_1,\ldots,x_d)$ and consider
the standard one-hidden-layer tanh network
\begin{equation}\label{eq:network}
 z_{W,v}(x)=\sum_{a=1}^N v_a\tanh\!\left(\sum_{i=0}^d W_{ai}\bar x_i\right).
\end{equation}
There are no biases other than the constant input coordinate.
Choose a number $\delta\in(0,1)$ and set
\begin{equation}\label{eq:lambda}
 \lambda_0=1,\qquad \lambda_i=1+\delta 2^{-i}\quad(1\le i\le d).
\end{equation}
Independently, let
\begin{equation}\label{eq:prior}
 W_{ai}\sim N\!\left(0,\frac{N\sigma^2}{d\lambda_i}\right),
 \qquad v_a\sim N\!\left(0,\frac{\sigma^2}{N}\right).
\end{equation}
Its density is proportional to $\exp[-\kappa_{N,\lambda}/(2\sigma^2)]$,
where
\begin{equation}\label{eq:energy}
 \kappa_{N,\lambda}(W,v)
 =\frac dN\sum_{i=0}^d\lambda_i\|W_{:,i}\|_2^2+N\|v\|_2^2.
\end{equation}
The parameter $\sigma$ is a common prior scale. Multiplying every variance
in~\eqref{eq:prior} by the same positive width-dependent factor merely
reparametrizes this scale at a fixed width.

Before observing labels, condition the prior on
\begin{equation}\label{eq:global-margin}
 |z_{W,v}(x)|>\frac12\qquad\text{for every }x\in\{-1,1\}^d.
\end{equation}
After each prediction, reveal the true label and condition the posterior
on its exact sign. The strong-prior limit is $\sigma\downarrow0$ at fixed
$d,N,\delta$. No width limit is taken in the proof.

For any desired $\eta>0$, choosing $0<\delta<\min\{1,\eta\}$ gives
$1\le\lambda_i\le1+\eta$ for every coordinate, uniformly in $d$ and $N$.
The activation remains standard tanh. The model change is precisely the
small coordinate dependence in~\eqref{eq:prior}; it is not the exactly
isotropic first-layer law.

\begin{theorem}\label{thm:main}
Let $d\ge17$ be odd, put $k=\lfloor d/8\rfloor$ and $q=2k-1$, and use
the target
\[
 f(x)=\sgn\!\left(\sum_{i=1}^d x_i\right).
\]
There is a fixed permutation of the input cube, depending only on $d$,
with a designated set of $T$ rounds such that
\begin{equation}\label{eq:mistake-count}
 T\ge\frac{\binom d k}{\binom{2k-1}k}\ge2^k.
\end{equation}
For every $N\ge1$ and every $\delta\in(0,1)$, the posterior probability
of the correct label on each designated round tends to zero as
$\sigma\downarrow0$ under~\eqref{eq:prior}--\eqref{eq:global-margin}.

Consequently, for every $\alpha\in(0,1)$ and every fixed $N,\delta$,
all sufficiently small positive $\sigma$ have the following properties
on that one fixed permutation:
\begin{enumerate}
\item Fresh posterior sampling makes all $T$ designated mistakes with
probability at least $1-\alpha$.
\item Posterior-label majority is wrong on all $T$ designated rounds.
\item The strict $2/3$ confidence requirement fails on all $T$
designated rounds.
\end{enumerate}
The target has bounded-weight teacher complexity $C_2(f)=1$ and Boolean
circuit size $\Theta(d)$ for a fixed bounded-fan-in complete gate basis.
Thus every learner width $N\ge1$ satisfies $N\ge C_2(f)$.
\end{theorem}

The variance threshold may depend on $N,d,\delta,\alpha$.
The theorem does not assert a positive threshold for this $\sigma$
uniform over all widths. Proposition~\ref{prop:uniform-scale} gives
a width-uniform threshold in the mean-field scale $t=\sqrt N\sigma$.
The covariance coefficients and input permutation themselves are the
same for every width. The calibration prefix contains all inputs in two
extreme Hamming regions; no short-calibration claim is made.

\section{An energy bound for every nonlinear classifier}

A coordinate $i\in\{1,\ldots,d\}$ is \emph{essential} to a Boolean
function $g$ if some pair of inputs differing only in coordinate $i$
receives opposite labels. Let $\ess(g)$ be its set of essential input
coordinates. A function depends only on this set: changing any other
coordinate never changes its value.

Define
\begin{equation}\label{eq:phi}
 \Phi_d(Q)=\min_{u>0}\left\{dQu^2+\frac14\coth^2u\right\},
 \qquad Q>0.
\end{equation}
For a represented function $g$, let $E_{N,\lambda}(g)$ be the infimum
of~\eqref{eq:energy} over parameters with
$g(x)z_{W,v}(x)>1/2$ at every cube input.

\begin{lemma}[Pair contraction]\label{lem:pair}
For real $c,t$,
\[
 \frac{|\tanh(c+t)-\tanh(c-t)|}{2}\le\tanh|t|.
\]
The function $s\mapsto\tanh^2\sqrt{s}$ is concave on $[0,\infty)$.
Consequently, for $w\in\R^N$,
\begin{equation}\label{eq:jensen}
 \sum_{a=1}^N\tanh^2|w_a|
 \le N\tanh^2\!\left(\frac{\|w\|_2}{\sqrt N}\right).
\end{equation}
\end{lemma}
\begin{proof}
The identity
\[
 \tanh(c+t)-\tanh(c-t)
 =\frac{\sinh(2t)}{\cosh(c+t)\cosh(c-t)}
\]
and
$\cosh(c+t)\cosh(c-t)=(\cosh(2c)+\cosh(2t))/2$
show that the absolute difference is largest at $c=0$.
For concavity, put $t=\sqrt{s}>0$. The derivative in $s$ is
$\tanh(t)\operatorname{sech}^2(t)/t$; its derivative in $t$ has the sign of
\[
 t(1-3\tanh^2t)-\tanh t<0.
\]
Indeed, $\tanh t\ge t\operatorname{sech}^2t$, because
$\sinh(2t)\ge2t$, and substitution bounds the displayed expression
by $-2t\tanh^2t<0$. Continuity handles $s=0$.
Jensen's inequality gives~\eqref{eq:jensen}.
\end{proof}

\begin{lemma}[Weighted essential coordinates]\label{lem:essential}
For every nonconstant represented Boolean function $g$ and every width
$N\ge1$,
\begin{equation}\label{eq:essential-lower}
 E_{N,\lambda}(g)\ge
 \Phi_d\!\left(\sum_{i\in\ess(g)}\lambda_i\right).
\end{equation}
If $S\subseteq\{1,\ldots,d\}$ has odd cardinality and
$h_S(x)=\sgn(\sum_{i\in S}x_i)$, then equality holds:
\begin{equation}\label{eq:majority-energy}
 E_{N,\lambda}(h_S)=\Phi_d\!\left(\sum_{i\in S}\lambda_i\right).
\end{equation}
The function $\Phi_d$ is strictly increasing.
\end{lemma}
\begin{proof}
It suffices to work with weak signed margins at least $1/2$.
The weak and strict infima agree: multiply the readout of a weakly
feasible network by $1+\rho$ and let $\rho\downarrow0$.

Fix weakly feasible parameters and put $C=\sqrt N\|v\|_2$.
For each essential $i$, choose adjacent inputs $x,x'$ with opposite
labels. Their output difference has magnitude at least one. Writing
the hidden preactivations as $c_a\pm W_{ai}$, applying
Lemma~\ref{lem:pair}, and using Cauchy--Schwarz gives
\begin{align*}
 \frac12
 &\le\frac{|z(x)-z(x')|}{2}\\
 &\le\|v\|_2\left(\sum_{a=1}^N\tanh^2|W_{ai}|\right)^{1/2}
 \le C\tanh\!\left(\frac{\|W_{:,i}\|_2}{\sqrt N}\right).
\end{align*}
In particular $C>1/2$. Set $u=\operatorname{arctanh}(1/(2C))>0$.
Every essential column satisfies $\|W_{:,i}\|_2^2/N\ge u^2$.
The bias and other columns contribute nonnegative energy, so
\[
 \kappa_{N,\lambda}(W,v)
 \ge d\left(\sum_{i\in\ess(g)}\lambda_i\right)u^2+C^2
 =d\left(\sum_{i\in\ess(g)}\lambda_i\right)u^2+\frac14\coth^2u.
\]
Minimizing proves~\eqref{eq:essential-lower}.

For odd $|S|$, set every row of $W$ equal to $u$ on coordinates in
$S$ and zero elsewhere, and set every $v_a=\coth(u)/(2N)$.
The output is
\[
 z(x)=\frac12\coth(u)\tanh\!\left(u\sum_{i\in S}x_i\right).
\]
The integer sum has absolute value at least one, so the network has
majority sign and weak margin at least $1/2$. Its energy is exactly
the objective in~\eqref{eq:phi} with $Q=\sum_{i\in S}\lambda_i$.
Readout scaling and the lower bound prove~\eqref{eq:majority-energy}.

For $Q>0$, the objective in~\eqref{eq:phi} tends to infinity at both
endpoints $u\downarrow0$ and $u\to\infty$, and therefore has a
minimizer $u_Q>0$. If $Q_2>Q_1$, then
\[
 \Phi_d(Q_2)
 \ge\Phi_d(Q_1)+d(Q_2-Q_1)u_{Q_2}^2>\Phi_d(Q_1).
\]
\end{proof}

The lower bound applies to every wide nonlinear network. In particular,
it supplies a comparison with functions outside the majority family,
which cannot be obtained by examining scalar candidate subnetworks alone.

\section{Calibration and an exact minimum-energy tier}

Write $\wt(x)=|\{i:x_i=+1\}|$ and put $q=2k-1$.
Reveal every input satisfying $\wt(x)\le k-1$ with label $-1$,
and every input satisfying $\wt(x)\ge d-k+1$ with label $+1$.
Use any fixed order within this prefix. All its labels agree with
the full majority target because $k<d/2$.

\begin{lemma}[Two-sided calibration]\label{lem:calibration}
Every Boolean function consistent with the calibration prefix has at
least $q$ essential input coordinates. If it has exactly $q$, then it
is $h_S$ for its essential support $S$, with $|S|=q$.
Every such majority $h_S$ is consistent with the prefix.
\end{lemma}
\begin{proof}
Suppose a consistent function depends on a support of size
$r\le2k-2$. Choose an assignment on that support with $t$ positive
coordinates, where
\[
 \max\{0,r-k+1\}\le t\le\min\{r,k-1\}.
\]
This interval is nonempty. Setting all outside coordinates negative
gives a negative calibration input, while setting all of them positive
gives a positive calibration input, since $t+d-r\ge d-k+1$.
The two inputs agree on the support, a contradiction.

If $r=q=2k-1$, any support assignment with at most $k-1$ positives
extends to a negative calibration input by setting outside bits
negative. Any assignment with at least $k$ positives extends to a
positive calibration input by setting outside bits positive, because
$k+d-q=d-k+1$. These cases exhaust the support cube and force
majority. Conversely, a $q$-coordinate majority is negative when
the full input has at most $k-1$ positives, and positive when the
full input has at most $k-1$ negatives.
\end{proof}

\begin{proposition}[Uniform separation of all competitors]\label{prop:tier}
At every width $N$, all functions $h_S$ with $|S|=q$ have pairwise
distinct energies, and each is strictly cheaper than every other
initially consistent Boolean function. After any further labels,
while at least one of these candidates survives, the unique
minimum-energy consistent function is the surviving $h_S$ with
smallest $\sum_{i\in S}2^{-i}$.
\end{proposition}
\begin{proof}
By Lemma~\ref{lem:essential}, the candidate energy is exactly
\[
 \Phi_d(Q_S),\qquad
 Q_S=\sum_{i\in S}\lambda_i=q+\delta\sum_{i\in S}2^{-i}.
\]
Finite binary subset sums are distinct. For example, at the smallest
index belonging to exactly one of two subsets, its contribution
$2^{-i}$ is larger than the sum of all later contributions through
index $d$. Strict monotonicity of $\Phi_d$ therefore separates the
candidate energies. Also
\begin{equation}\label{eq:gap}
 Q_S<q+\delta<q+1.
\end{equation}
Every outsider has at least $q+1$ essential coordinates by
Lemma~\ref{lem:calibration}, and every $\lambda_i\ge1$.
Its energy is at least $\Phi_d(q+1)$, strictly above every candidate
energy by~\eqref{eq:gap}. The same comparison holds after further
conditioning whenever a candidate remains.
\end{proof}

Both the separation from outsiders and the ordering within the candidate
family are independent of $N$. The ordering within the family is also
independent of the positive value of $\delta$.

\section{The input order and posterior mistakes}

After calibration, repeat the following deterministic procedure while
some candidate remains. Let $S$ be the surviving $q$-set with smallest
$\sum_{i\in S}2^{-i}$. Choose its lexicographically first $k$-element
subset $A$ and query $x^A$, the vector positive exactly on $A$.
Its true majority label is negative because $k<d/2$.
The selected candidate satisfies $h_S(x^A)=+1$.

This query is new: any previously queried negative input would have
eliminated a candidate that predicted positively there. It is outside
the calibration prefix because it has weight exactly $k$.
It eliminates exactly those candidate supports containing $A$.
Indeed, a majority on $2k-1$ coordinates predicts positively at a
weight-$k$ input precisely when all its $k$ positive coordinates lie
in that support.

There are initially $\binom d{2k-1}$ candidates. Each query eliminates
at most $\binom{d-k}{k-1}$ of them, so exhausting the candidate class
requires at least
\[
 T\ge\frac{\binom d{2k-1}}{\binom{d-k}{k-1}}
 =\frac{\binom d k}{\binom{2k-1}k}
 \ge\frac{(d/k)^k}{4^k}\ge2^k.
\]
Here $\binom d k\ge(d/k)^k$, $\binom{2k-1}k\le4^k$, and $d\ge8k$.
The entire sequence is computable before the run. It depends only on
the candidate supports and the fixed negative labels, not on posterior
samples. Complete it to a permutation by revealing all remaining cube
inputs in a fixed order with their majority labels.

\begin{lemma}[Gaussian concentration on a unique minimum]\label{lem:rate}
At a fixed finite width and after any fixed feasible label history,
if a unique consistent Boolean function has the least value of
$E_{N,\lambda}$, then its posterior probability tends to one as
$\sigma\downarrow0$.
\end{lemma}
\begin{proof}
For a represented Boolean function $g$, write
\[
 U_g=\{(W,v):g(x)z_{W,v}(x)>1/2\text{ for every }x\}.
\]
This is open. Its energy infimum agrees with that over its closed
weak-margin set, by readout scaling. The latter set is nonempty and
the positive definite quadratic energy is coercive, so the weak
infimum is attained. In the finite-dimensional Gaussian parameter
space,
\begin{equation}\label{eq:ldp}
 \lim_{\sigma\downarrow0}\sigma^2\log\Pp_\sigma(U_g)
 =-\frac12 E_{N,\lambda}(g).
\end{equation}
For an elementary lower bound, take a small ball contained in $U_g$
whose maximum energy is arbitrarily close to its infimum; Gaussian
normalizing powers of $\sigma$ disappear after multiplication by
$\sigma^2$. For the upper bound, write the energy as
$(1-\varepsilon)\kappa+\varepsilon\kappa$, bound the first part below
by $(1-\varepsilon)E_{N,\lambda}(g)$, and integrate the second Gaussian
factor over the entire parameter space. The resulting constant
$\varepsilon^{-P/2}$, where $P$ is the parameter dimension, also
disappears at the logarithmic scale. Then let $\varepsilon\downarrow0$.

There are finitely many Boolean functions on the cube.
The globally margin-conditioned posterior after a label history is
the prior restricted to the union of the compatible events $U_g$.
A unique minimum energy has a positive gap from all other represented
compatible functions. Equation~\eqref{eq:ldp} makes the ratio of each
other event's probability to that of the minimum tend to zero.
\end{proof}

By Proposition~\ref{prop:tier}, every designated query is predicted
positively by the unique minimum-energy consistent function, while its
true label is negative. Lemma~\ref{lem:rate} proves the claimed
correct-label probability limit. For fixed $N,\delta$ there are only
finitely many designated rounds, so one sufficiently small positive
variance makes the correct-label probability smaller than
$\min\{\alpha/T,1/3\}$ on each. A union bound shows that all designated
posterior samples are wrong with probability at least $1-\alpha$.
The majority and confidence conclusions follow from the same bound.

\paragraph{Target resources.}
The width-one teacher $z(x)=\tanh(\sum_i x_i)$ has every weight at
most one, and its signed margin is at least $\tanh(1)>1/2$ because
$d$ is odd. Thus $C_2(f)=1$. For a linear-size Boolean circuit,
convert the input bits to zero-one form and count them with carry-save
adders. Each constant-size full adder replaces three equally weighted
bits by two, decreasing the number of stored bits by one. At most
$O(d)$ such adders leave two $O(\log d)$-bit numbers; their final
addition and comparison with $d/2$ use $O(\log d)$ further gates.
This gives an $O(d)$-gate circuit. Every input is essential to majority,
so bounded fan-in gives an $\Omega(d)$ gate lower bound. This completes
the proof of Theorem~\ref{thm:main}.

\section{A linear calibration prefix}

The preceding result permits every fixed $\delta\in(0,1)$, but its
calibration is large. A compactness argument gives a shorter prefix
when the perturbation is sufficiently small for each dimension.
The quantitative bound below is explicit; we first give the shorter
compactness argument explaining why a width-independent gap exists.

\begin{theorem}[Linear calibration]\label{thm:short}
Let $q\ge3$ be odd and $d=9q$. For every
\begin{equation}\label{eq:explicit-delta}
 0<\delta\le\frac{1}{80000q^3},
\end{equation}
the model
\eqref{eq:network}--\eqref{eq:global-margin}, with the same full-majority
target, has a fixed input permutation beginning with exactly $2q$
calibration inputs and followed by exactly
\begin{equation}\label{eq:short-count}
 T=9^{(q+1)/2}=\exp(\Omega(d))
\end{equation}
designated mistake rounds. For every width $N\ge1$, the correct-label
posterior probability on each designated round tends to zero as
$\sigma\downarrow0$. All finite-positive-variance conclusions of
Theorem~\ref{thm:main} follow as well.

The chosen $\delta$ and input permutation are independent of $N$.
The perturbation can consequently be arbitrarily small. Its explicit
upper bound depends on $d$; no positive dimension-uniform amplitude
is asserted in this linear-calibration theorem.
\end{theorem}

\subsection{A uniform gap above the majority minimum}

The following relaxation is used only to prove an energy gap uniform
over all widths. It does not replace the finite-width Gaussian posterior.
For a probability measure $\nu$ on $\R^{d+1}\times\R$, write its
variables as $(a,b)$ and define
\begin{equation}\label{eq:measure-model}
 z_\nu(x)=\int b\tanh(a\cdot\bar x)\,d\nu(a,b),\qquad
 I(\nu)=\int(d\|a\|_2^2+b^2)\,d\nu(a,b).
\end{equation}
Every width-$N$ network is included: take
$\nu=N^{-1}\sum_{j=1}^N\delta_{(W_{j,:},Nv_j)}$.
Then $z_\nu=z_{W,v}$ and $I(\nu)=\kappa_{N,\boldsymbol 1}(W,v)$.

\begin{lemma}[Compactness and equality]\label{lem:compact-gap}
For fixed $d$ and odd $q\le d$, there is $\Delta=\Delta(d,q)>0$
such that every globally weak-margin classifier represented
by~\eqref{eq:measure-model} and having at least $q$ essential input
coordinates either is a signed majority on exactly $q$ coordinates
or has energy at least $\Phi_d(q)+\Delta$.
\end{lemma}
\begin{proof}
First, energy sublevels are sequentially compact for the weak topology
of probability measures. The second-moment bound gives tightness, and
the integral of the nonnegative continuous quadratic energy is lower
semicontinuous. For each fixed cube input, $z_\nu(x)$ is continuous
along an energy-bounded weakly convergent sequence. Indeed,
$|b\tanh(a\cdot\bar x)|\le|b|$ and
\[
 \int |b|\mathbf1_{\{|b|>R\}}\,d\nu\le\frac{I(\nu)}R.
\]
Truncating the $b$ coordinate gives bounded continuous integrands and
the displayed uniform tail bound removes the truncation. Consequently,
for any fixed Boolean function, its minimum weak-margin energy over
measures is attained whenever its feasible set contains a measure of
finite energy. This uses only tightness of probability measures on a
finite-dimensional Euclidean space and the finite number of margin
constraints.

The essential-coordinate proof of Lemma~\ref{lem:essential} applies
with integrals in place of empirical averages. If the essential support
has size $r$, it gives $I(\nu)\ge\Phi_d(r)$. We examine equality for
$r=q$ at $I(\nu)=\Phi_d(q)$.

Put $C=(\int b^2\,d\nu)^{1/2}$ and
$u=\operatorname{arctanh}(1/(2C))>0$. Equality in the summed energy
bound forces $\int a_j^2\,d\nu=u^2$ for every essential $j$, and
all unused coordinates, including $a_0$, to vanish almost surely.
For each essential $j$, fix an adjacent pair with opposite labels.
All inequalities in the pair-contraction, Jensen, and Cauchy--Schwarz
chain must be equalities; otherwise its column-norm bound is strict
and the total energy exceeds $\Phi_d(q)$.

Strict concavity in Lemma~\ref{lem:pair} gives $|a_j|=u$ almost surely.
Equality in the scalar pair contraction forces the center of the two
preactivations to be zero, because $u>0$.
Equality in Cauchy--Schwarz then gives, for a sign
$\gamma_j\in\{-1,1\}$ fixed by the orientation of the chosen pair,
\[
 b=\gamma_j C\sgn(a_j)\quad\text{almost surely}.
\]
Thus $|b|=C$ and $a_j=(u/C)\gamma_j b$ for every essential $j$.
Intersecting these finitely many full-measure events yields
\[
 z_\nu(x)
 =C\tanh\!\left(u\sum_{j\in S}\gamma_jx_j\right).
\]
Its Boolean sign is a signed $q$-coordinate majority. This proves
that no other function with exactly $q$ essential coordinates can
attain energy $\Phi_d(q)$.

There are finitely many Boolean functions on the cube. Attainment
therefore gives a common positive gap to all finite-energy,
non-majority functions with exactly $q$ essential coordinates.
Functions with more than $q$ essential coordinates have the positive
gap $\Phi_d(q+1)-\Phi_d(q)$. The minimum of these finitely many
positive gaps is a suitable $\Delta$. Functions with no finite-energy
representation need not be included in the minimum.
\end{proof}

\subsection{Block calibration and elimination}

Partition the $d=9q$ coordinates into $q$ blocks $B_1,\ldots,B_q$
of size nine, with block indices read cyclically. Put $m=(q-1)/2$.
For each $j$, use two block-constant calibration inputs. In the
positive one, the positive blocks are
$\{j,j+1,\ldots,j+m\}$; in the negative one, remove block $j$ from
this set. The inputs in each pair differ only on block $B_j$, and
the other $q-1$ blocks are balanced. The resulting $2q$ inputs are
distinct and have the indicated full-majority labels.

A consistent Boolean function must have an essential coordinate in
every block, since otherwise it would give equal labels to that
block's two calibration inputs. In particular it has at least $q$
essential coordinates. A signed majority on exactly $q$ coordinates
must use one coordinate per block. Its sign on block $j$ must be
positive: a negative coefficient would make the output score
decrease when that block changes from the negative calibration
input to the positive one, contradicting the prescribed labels.
Conversely, every positive majority with one coordinate in each
block fits the calibration. Let $\mathcal C$ be these $9^q$ candidates.

Choose $\Delta$ from Lemma~\ref{lem:compact-gap}. If $u_q>0$ minimizes
\eqref{eq:phi} at $Q=q$, then
\[
 \Phi_d(q+\delta)\le\Phi_d(q)+d\delta u_q^2.
\]
For the qualitative argument, choose
$0<\delta<\min\{1,\Delta/(2du_q^2)\}$.
Appendix~\ref{sec:quantitative} supplies
$\Delta_q=1/(10000q^3)$ and $du_q^2\le4$, which justify every
$\delta$ in~\eqref{eq:explicit-delta}: candidate energies are then
strictly below $\Phi_d(q)+\Delta_q/2$, while outsiders have energy
at least $\Phi_d(q)+\Delta_q$. For the explicit conclusion, take
$\Delta=\Delta_q$ in the following comparison.
Every candidate has exact anisotropic energy
$\Phi_d(q+\delta\sum_{i\in S}2^{-i})<\Phi_d(q+\delta)$.
Every initially consistent outsider has isotropic energy at least
$\Phi_d(q)+\Delta$, and its anisotropic energy is no smaller because
$\lambda_i\ge1$. Thus every candidate is cheaper than every outsider,
uniformly over all finite widths. Distinct binary subset sums give
a unique cheapest surviving candidate at every stage.

Put $\ell=(q+1)/2$ and fix the first $\ell$ blocks. If $S$ is
the current cheapest surviving support, query the input positive only
at its chosen coordinates in those fixed blocks:
\[
 A=S\cap(B_1\cup\cdots\cup B_\ell),\qquad x=x^A.
\]
The query has $\ell$ positive coordinates, so its full-majority target
label is negative, while the selected candidate predicts positively.
A candidate $h_{S'}$ is positive there exactly when $S'$ contains
all $\ell$ coordinates of $A$: it needs at least $\ell$ positive
coordinates to win its $q$-coordinate majority, and the query has
only $\ell$ positives in total.

Thus each query deletes precisely one complete projection fiber,
consisting of the $9^{q-\ell}$ supports with that choice of one
coordinate in each of the fixed $\ell$ blocks. There are $9^\ell$
such fibers. Each selected candidate belongs to a surviving fiber,
so the procedure visits every fiber once and takes exactly
$T=9^\ell$ forced queries. This proves~\eqref{eq:short-count}.
No query repeats, and these queries lie outside the calibration
prefix: their weight is $\ell$, whereas calibration weights are
$9(q-1)/2$ and $9(q+1)/2$.

The procedure uses only the fixed candidate order and previous target
labels, so it defines one fixed input sequence, completed arbitrarily
to a permutation. The Gaussian concentration and finite-variance
arguments from Lemma~\ref{lem:rate} apply at every fixed width without
change. This proves Theorem~\ref{thm:short}.

\section{A common mean-field scale at every width}

The preceding theorems use the scale $\sigma$ in~\eqref{eq:prior}.
There is also a natural parametrization with a single positive scale
that works uniformly over all widths. Put $t=\sqrt N\sigma$, so that
\begin{equation}\label{eq:meanfield-prior}
 W_{ai}\sim N\!\left(0,\frac{t^2}{d\lambda_i}\right),\qquad
 v_a\sim N\!\left(0,\frac{t^2}{N^2}\right).
\end{equation}
This is the same family of priors with a width-dependent change of
scale variable.

\begin{proposition}[Uniform small mean-field scale]\label{prop:uniform-scale}
Fix either construction above, including its allowed positive
anisotropy and its designated $T$ rounds. For every $\alpha\in(0,1)$
there is $t_0>0$, independent of $N$, such that for every $N\ge1$
and every $0<t<t_0$, the posterior under~\eqref{eq:meanfield-prior}
makes all $T$ designated sampling mistakes with probability at least
$1-\alpha$. The posterior-majority and confidence conclusions are
uniform in $N$ as well.
\end{proposition}
\begin{proof}
At a designated round, let $h_S$ be the unique minimum-energy
consistent function and let $E=\Phi_d(Q_S)$. The proofs above give a
number $g>0$, independent of $N$, such that every consistent function
predicting the next label correctly has energy at least $E+g$.
For the large-prefix construction this follows from the finitely
many explicit candidate gaps and the outsider bound $\Phi_d(q+1)$.
For the short-prefix construction use the uniform outsider gap from
Lemma~\ref{lem:compact-gap} and the same explicit candidate gaps.

Write $a_a=W_{a,:}$ and $b_a=Nv_a$. Under~\eqref{eq:meanfield-prior},
the pairs $(a_a,b_a)$ are IID Gaussian, with coordinate variances
$t^2/(d\lambda_i)$ and $t^2$. Let
\[
 J(a,b)=d\sum_{i=0}^d\lambda_i a_i^2+b^2,\qquad P=d+2.
\]
Then $\kappa_{N,\lambda}=N^{-1}\sum_aJ(a_a,b_a)$ and
$J(a_a,b_a)/t^2$ is chi-squared with $P$ degrees of freedom.

A cloned majority minimizer is a single pair $(a,b)$ of cost $E$.
Slightly enlarge its readout and take a sufficiently small open ball
$B$ about the resulting pair. Every pair in $B$ has the same strict
signed margins as $h_S$, and $\sup_B J\le E+g/4$.
If all $N$ pairs lie in $B$, their averaged output has those strict
margins and fits the current history. Its prior probability is
$p_t^N$, where Gaussian density integration gives
\[
 p_t\ge c_B t^{-P}\exp\!\left(-\frac{E+g/4}{2t^2}\right)
\]
for a positive constant $c_B$ independent of $N,t$.

Set $\zeta=g/[4(E+g+1)]\in(0,1)$. The chi-squared moment-generating
function gives
\[
 \Pp_t\{\kappa_{N,\lambda}\ge E+g\}
 \le\left[\zeta^{-P/2}
       \exp\!\left(-\frac{(1-\zeta)(E+g)}{2t^2}\right)\right]^N.
\]
The numerator of the posterior correct-label probability is bounded
by this tail; its denominator is at least $p_t^N$. Since
$\zeta(E+g)\le g/4$, their ratio is at most
\[
 \left[c_B^{-1}\zeta^{-P/2}t^P
              \exp\!\left(-\frac{g}{4t^2}\right)\right]^N
 =r(t)^N.
\]
Here $r(t)\to0$ as $t\downarrow0$. Once $r(t)<1$, this bound is at
most $r(t)$ for every $N\ge1$. There are finitely many designated
rounds, so choose a common $t_0$ making every bound smaller than
$\min\{\alpha/T,1/3\}$. A union bound finishes the proof.
\end{proof}

\section{Polynomial IID learning for the same prior}

The online lower bounds do not arise from a lack of polynomial IID
sample guarantees for this target. The following bound is independent
of the learner width and applies to both constructions.

\begin{theorem}[IID upper bound]\label{thm:iid}
Let $d$ be odd, $0<\delta<1$, and use the majority target and prior
\eqref{eq:lambda}--\eqref{eq:global-margin}. Define
\[
 K=5d+1,\qquad
 m=\left\lceil32K\bigl((d+1)\log2+1\bigr)\right\rceil,
 \qquad B_d=m\bigl((d+1)^2+1\bigr)\log2.
\]
For every width $N\ge1$, every input distribution $\mu$, and every
$\rho,\beta\in(0,1)$, if
\begin{equation}\label{eq:iid-n}
 n\ge\frac6\rho\left(B_d+\log\frac1\beta\right),
\end{equation}
then, at all sufficiently small positive prior scales, with probability
at least $1-\beta$ over $n$ IID labeled examples, the $\mu$-mass of
inputs whose posterior correct-label probability is at most $2/3$
is at most $\rho$.
The sample bound is $O((d^4+\log(1/\beta))/\rho)$.
For fixed $N,d,\delta,\rho$, the sufficiently-small-scale threshold
can be chosen uniformly over all target-consistent histories and
sample sizes. In the mean-field parametrization~\eqref{eq:meanfield-prior},
one sufficiently small positive $t$ works uniformly over all widths
as well.
\end{theorem}
\begin{proof}
We first give an energy-bounded realization of the target. Put
$a=d^{-1/2}$, give every hidden row coefficient $a$ on every input
coordinate and zero bias, and give every readout coefficient
$1/(N\tanh a)$. Its output is
\[
 z(x)=\frac{\tanh(a\sum_i x_i)}{\tanh a},
\]
which has signed margin at least one. Its energy is
\[
 K_0=\sum_{i=1}^d\lambda_i+\coth^2a<d+1+4d=K,
\]
using $\sum_i\lambda_i<d+1$ and $\tanh a\ge a/2$ for $0<a\le1$.

Next, every width-$N$ network with energy at most $K$ satisfies
$\sum_a|v_a|\le\sqrt N\|v\|_2\le\sqrt K$.
If $U$ is uniform on $(-1,1)$, then
\[
 \tanh t=\E\,\sgn\bigl(t-\operatorname{arctanh}U\bigr).
\]
Choose neuron $a$ with probability $|v_a|/\sqrt K$, draw $U$, and
output the signed affine threshold
\[
 h(x)=\sgn(v_a)\sgn\bigl(W_{a,:}\cdot\bar x-
                     \operatorname{arctanh}U\bigr).
\]
Use the zero function with the remaining probability. Then
$h(x)\in\{-1,0,1\}$ and $\E h(x)=z(x)/\sqrt K$.
For $m$ independent voters, Hoeffding's inequality and a union bound
over the cube give
\[
 \Pp\!\left\{\max_x\left|
       \frac{\sqrt K}{m}\sum_{j=1}^m h_j(x)-z(x)\right|\ge\frac14\right\}
 \le2^{d+1}e^{-m/(32K)}<1.
\]
Thus a globally weak-margin network of energy at most $K$ has the
same sign as some sum of $m$ such voters. Almost surely each nonzero
voter is a strict affine threshold on the finite cube, and its negative
is again a strict affine threshold. The approximating sum has no zeros
because the network's absolute margin is at least $1/2$.

There are at most
\[
 H_d\le2\sum_{j=0}^d\binom{2^d-1}j\le2^{(d+1)^2}
\]
strict affine threshold functions on the cube, by the central
hyperplane arrangement bound in coefficient dimension $d+1$.
Allowing the zero voter, the class $\mathcal H_K$ of all resulting
sign functions has size at most
$(H_d+1)^m\le e^{B_d}$. Every globally margin-separated
energy-$K$ function belongs to this one finite class, uniformly over
all learner widths.

The posterior mass of parameters with energy greater than $K$ tends
to zero uniformly over target-consistent histories. To see this, let
$U_f$ be the target's strict global-margin event. It contains an open
ball of energy strictly below $K$, by the witness above.
Every target-consistent posterior conditioning event contains $U_f$,
so, using probabilities under the unconditioned Gaussian prior,
\[
 \Pp_\sigma\{\kappa>K\mid\text{global margin and history}\}
 \le\frac{\Pp_\sigma\{\kappa>K\}}{\Pp_\sigma(U_f)}.
\]
The Gaussian tail in the numerator has rate at least $K/2$ and the
denominator has a strictly smaller rate, by the open ball. Hence the
ratio tends to zero at each fixed width. It is independent of the
history and sample size. Choose the scale small enough to make it
at most $\rho/6$. In the mean-field parametrization, use the
product-ball and chi-squared estimate from
Proposition~\ref{prop:uniform-scale} with $E=K_0$ and $g=K-K_0$;
the strict target witness supplies the denominator ball. This makes
the same tail bound uniform over every width, in addition to every
history.

For a fixed $g\in\mathcal H_K$ of $\mu$-error greater than $\rho/6$,
the probability that it fits $n$ IID examples is at most
$e^{-n\rho/6}$. A union bound and~\eqref{eq:iid-n} show that, with
probability at least $1-\beta$, every consistent function in
$\mathcal H_K$ has error at most $\rho/6$. On this event, the
posterior-averaged test error is at most $\rho/6+\rho/6=\rho/3$.
At any input whose correct-label probability is at most $2/3$, the
posterior wrong-label probability is at least $1/3$. Integrating
therefore bounds the mass of such inputs by $\rho$.
\end{proof}

\begin{corollary}[Fixed exponents above one]\label{cor:gamma}
Fix either online construction, including its permitted anisotropy.
Consider the independent-Gaussian prior
\begin{equation}\label{eq:gamma-prior}
 \operatorname{Var}(W_{ai})=\frac{\tau^2N^{1-\gamma}}{d\lambda_i},
 \qquad
 \operatorname{Var}(v_a)=\tau^2N^{-\gamma-1},
\end{equation}
with the same margin and exact-label conditioning.
\begin{enumerate}
\item For every fixed $\gamma>1$, the unit-prefactor choice $\tau=1$
has all the online conclusions at every sufficiently large width.
The polynomial IID bound in Theorem~\ref{thm:iid} holds there as well.
\item For every prescribed online failure tolerance and IID accuracy
parameters, one sufficiently small positive $\tau$ makes both
conclusions hold for every $N\ge1$ and every $\gamma\ge1$
simultaneously.
\end{enumerate}
The covariance asymmetry remains present. No conclusion here covers
$\gamma=1$ with unit prefactor $\tau=1$.
\end{corollary}
\begin{proof}
The mean-field scale is $t_N=\tau N^{(1-\gamma)/2}$.
For fixed $\gamma>1$ and $\tau=1$, it tends to zero as $N\to\infty$.
For $\gamma\ge1$ and $N\ge1$, it is at most $\tau$.
Apply the width-uniform thresholds in
Proposition~\ref{prop:uniform-scale} and Theorem~\ref{thm:iid}.
\end{proof}

\section{What exact isotropy changes}

The strict energy ordering inside the majority tier is necessary for
this proof. It cannot be inferred by choosing one preferred minimizer
when the covariance is exactly isotropic.

\begin{proposition}[The isotropic tier]\label{prop:isotropic}
Set $\delta=0$ and keep the same two-sided calibration.
At every width, the minimum-energy consistent functions are exactly
the $q$-coordinate majorities. The strong-prior posterior is uniform
over whichever of them remain consistent. Under any continuation,
its expected number of sampling mistakes before this tier is exhausted
is at most
\[
 1+\log\binom d q.
\]
This bounds only this tier, not the entire later learning process.
\end{proposition}
\begin{proof}
Lemmas~\ref{lem:essential} and~\ref{lem:calibration} give the common
candidate energy $\Phi_d(q)$ and outsider lower bound $\Phi_d(q+1)$.
Coordinate permutations preserve the isotropic Gaussian law and the
global-margin event, and act transitively on $q$-coordinate majorities.
Thus their prior function probabilities are exactly equal at every
positive variance, including all parameter-volume factors.
Conditioning on exact labels preserves equal masses among the survivors.
The outsider energy gap implies that the cold posterior is uniform
on those survivors.

If a query eliminates a fraction $p<1$ of a surviving class of size
$m$, its sampling-error probability is $p$ and the decrease in its
logarithmic size is $-\log(1-p)\ge p$. Summing gives at most
$\log\binom d q$ before the final exhaustion, which costs at most
one additional expected mistake. The same conditional argument applies
to an adaptive continuation that does not observe the current posterior
sample before choosing its query.
\end{proof}

An arbitrarily small positive $\delta$ splits the candidate energies
and yields exponentially many certain cold-limit mistakes. At
$\delta=0$, the equal-mass mixture has the logarithmic bound above.
The limits of zero covariance asymmetry and zero prior scale therefore
need not commute. The theorem establishes the all-width lower bound
for the stated perturbed covariance and leaves the corresponding
exactly isotropic question open.

\appendix
\section{An explicit rigidity gap}\label{sec:quantitative}

We prove the numerical bound used in Theorem~\ref{thm:short}.
This argument applies both to finite networks and to the probability
measures in~\eqref{eq:measure-model}.

\begin{lemma}[Quantitative rigidity]\label{lem:quantitative}
Let $q\ge3$ be odd, $d=9q$, and
$\Delta_q=1/(10000q^3)$. If a globally weak-margin Boolean function
has exactly $q$ essential input coordinates and a representation
of isotropic energy at most $\Phi_d(q)+\Delta_q$, it is a signed
majority on those coordinates. Every function with more than $q$
essential input coordinates has energy at least
$\Phi_d(q)+27/64$.
\end{lemma}
\begin{proof}
Write the representation as in~\eqref{eq:measure-model}, and let $S$
be its essential support of size $q$. Put
\[
 K=I(\nu),\quad C=(\E_\nu b^2)^{1/2},\quad
 u=\operatorname{arctanh}(1/(2C)),\quad t_0=\tanh u=\frac1{2C},
 \quad e=K-\Phi_d(q).
\]
For each essential $j$, take an oppositely labeled adjacent pair,
ordered by coordinate $x_j=+1$ and $-1$. Let $s_j(a)$ be the
half-difference of its two hidden tanh outputs and let
$\varepsilon_j\in\{-1,1\}$ be its label orientation. Then
\[
 |s_j(a)|\le\tanh|a_j|,\qquad
 \sgn(s_j(a))=\sgn(a_j),\qquad
 \varepsilon_j\E_\nu[bs_j]\ge\frac12.
\]
With $r_j^2=\E_\nu a_j^2$, the essential-coordinate argument gives
$r_j^2\ge u^2$ and the nonnegative decomposition
\begin{align*}
 e={}&d\sum_{j\in S}(r_j^2-u^2)
       +d\E_\nu\sum_{i\notin S}a_i^2\\
    &+\bigl[dqu^2+C^2-\Phi_d(q)\bigr].
\end{align*}
The outside sum includes the bias coordinate. In particular,
\begin{equation}\label{eq:excess-columns}
 r_j^2-u^2\le e/d,\qquad
 \E_\nu\sum_{i\notin S}a_i^2\le e/d.
\end{equation}
Using the test value $u_*=1/\sqrt{6q}$ and
$\coth t\le1/t+t$ for $0<t\le1$ gives $\Phi_d(q)\le3q+1$.
Thus $e\le\Delta_q<1$ implies $K\le4q$, and
\begin{equation}\label{eq:u-range}
 C\le2\sqrt q,\qquad
 \frac1{4\sqrt q}\le u\le\frac2{3\sqrt q}<\frac12.
\end{equation}
The lower bound uses $u\ge1/(2C)$, and the upper one uses $dqu^2\le K$.

Put $h(s)=\tanh^2\sqrt s$. Besides $0<h'(s)\le1$, we have
\begin{equation}\label{eq:h-curvature}
 -h''(s)\ge\frac15\qquad(0\le s\le1).
\end{equation}
Indeed, writing $t=\sqrt s$ and $T=\tanh t$ gives
\[
 -h''(s)
 =\frac{\operatorname{sech}^2t\,[T-t+3tT^2]}{2t^3}
 \ge\operatorname{sech}^2t\,\frac{T^2}{t^2}.
\]
For $0<t\le1$, concavity of tanh and
$3/4<\tanh1<4/5$ imply that the last expression is greater than
$(9/25)(9/16)=81/400>1/5$; the endpoint follows by continuity.

For $u\le1/2$, define the tangent deficit
\[
 D_u(s)=h(u^2)+h'(u^2)(s-u^2)-h(s).
\]
For every real $a$,
\begin{equation}\label{eq:deficit}
 D_u(a^2)\ge\frac{u^2}{10}(|a|-u)^2.
\end{equation}
If $a^2\le1$,~\eqref{eq:h-curvature} gives
$D_u(a^2)\ge(a^2-u^2)^2/10$, which implies the claim.
If $a^2\ge1$, convexity of $D_u$ and~\eqref{eq:h-curvature} give
$D_u(1)\ge9/160$ and $D'_u(1)\ge3/20$.
Therefore $D_u(a^2)\ge(9/160)a^2$, which is stronger than
\eqref{eq:deficit} because $u^2\le1/4$.

Cauchy--Schwarz gives $\E s_j^2\ge t_0^2$, and
$\E h(a_j^2)\ge\E s_j^2$. Using~\eqref{eq:excess-columns} and
$h'\le1$ yields $\E D_u(a_j^2)\le e/d$, hence
\begin{equation}\label{eq:column-magnitude}
 \E(|a_j|-u)^2\le\frac{10e}{du^2}.
\end{equation}
Set $v=b/C$, so $\E v^2=1$. The oriented margin inequality gives
\begin{align}
 \|s_j-\varepsilon_jt_0v\|_{L^2(\nu)}^2
 &\le\E s_j^2-t_0^2\le h(r_j^2)-h(u^2)\le e/d,
 \label{eq:cs-stability}\\
 \|s_j-\tanh(a_j)\|_{L^2(\nu)}^2
 &\le\E[h(a_j^2)-s_j^2]\le e/d.
 \label{eq:pair-stability}
\end{align}
The second line uses the common signs and pointwise magnitude ordering
of $s_j$ and $\tanh(a_j)$.

Insert successively $u\sgn(a_j)$, $t_0\sgn(a_j)$,
$\tanh(a_j)$, and $s_j$, and use $u/t_0\le2$.
Equations~\eqref{eq:column-magnitude}--\eqref{eq:pair-stability} give
\begin{align*}
 \|a_j-\varepsilon_juv\|_{L^2(\nu)}
 &\le3\sqrt{\frac{10e}{du^2}}+4\sqrt{\frac ed}
 \le12\sqrt{\frac e{du^2}}\le16\sqrt e.
\end{align*}
Here~\eqref{eq:u-range} gives $du^2\ge9/16$.
The outside coordinates have $L^2$ contribution at most
$\sqrt{(d+1)e/d}\le2\sqrt e$ at every augmented cube input.
Consequently,
\[
 \left\|a\cdot\bar x-uv\sum_{j\in S}\varepsilon_jx_j\right\|_{L^2(\nu)}
 \le(16q+2)\sqrt e\le18q\sqrt e.
\]
By the Lipschitz property of tanh and Cauchy--Schwarz, the output
is uniformly within $36q^{3/2}\sqrt e$ of
\[
 F\!\left(\sum_{j\in S}\varepsilon_jx_j\right),\qquad
 F(t)=C\E\bigl[v\tanh(uvt)\bigr].
\]
For every nonzero $t$, the function $F(t)$ has the sign of $t$,
since its integrand has that sign whenever $v\ne0$ and $\E v^2=1$.
The signed sums are nonzero because $q$ is odd. If $e\le\Delta_q$,
the uniform approximation error is at most $36/100<1/2$.
A weakly margin-separated output cannot have the opposite sign
at any input. This proves the first assertion.

For the second assertion, Lemma~\ref{lem:essential} gives the lower
bound $\Phi_d(q+1)$. At its minimizer $u_{q+1}$,
\[
 \Phi_d(q+1)-\Phi_d(q)\ge d u_{q+1}^2,
 \qquad u_{q+1}^2\ge\frac1{4\Phi_d(q+1)}.
\]
The test value $u^2=1/(2\sqrt{d(q+1)})$ and the same coth bound give
$\Phi_d(q+1)\le\sqrt{d(q+1)}+1\le4(q+1)$.
Thus the gap is at least $9q/[16(q+1)]\ge27/64$.
\end{proof}

Finally, $du_q^2\le\Phi_d(q)/q\le4$ gives
\[
 \Phi_d(q+s)\le\Phi_d(q)+4s\qquad(s\ge0).
\]
Under~\eqref{eq:explicit-delta}, every candidate in the linear-prefix
construction has energy strictly below
$\Phi_d(q)+4\delta\le\Phi_d(q)+\Delta_q/2$.
Every outsider has energy at least $\Phi_d(q)+\Delta_q$, by
Lemma~\ref{lem:quantitative}; increasing coordinate weights cannot
reduce that energy. This completes the explicit quantitative
justification of Theorem~\ref{thm:short}.

\paragraph{Verification scope.}
The arguments are conventional mathematical proofs. They have undergone
independent AI-agent mathematical reviews and are not Lean formalized. The source, reproducible build, and local audit records are
included with this manuscript.

\begin{thebibliography}{1}
\bibitem{and-gaussian}
ChatGPT-6 Astra, Claude (Anthropic), Kaarel H\"anni, and Dmitry Vaintrob.
\emph{Online learning of AND under Gaussian neural-network priors}.
Research manuscript, 30 October 2026.
\href{../and-gaussian/paper.pdf}{Companion manuscript}, section on
an arbitrarily small anisotropy and its width-one majority corollary.
\end{thebibliography}

\end{document}
